% =====================================================================
%  05-transacciones.tex - Autorizacion, transacciones y recuperacion
% =====================================================================

\section{Autorización, transacciones y recuperación}\label{sec:pruebas-transacciones}

Esta sección verifica el modelo transaccional común de la \cref{sec:impl-http}: la autorización por recurso y su revocación, la atomicidad entre mutaciones y auditoría con el sellado concurrente, y la importación y restauración de la base. Las pruebas Rust de los adaptadores PostgreSQL y Redis usan el entorno aislado de la \cref{sec:pruebas-backends}; la demostración HTTP de \rutarepo{scripts/api-demo.sh}, que incorpora los ensayos de sellado concurrente y de migración, prepara sus propios servicios desechables. Las pruebas propias de cada dominio se presentan en sus secciones.

\subsection{Permisos por recurso y revocación}\label{sec:pruebas-permisos}

La demostración HTTP recorre cuatro roles y dos expedientes. Owner accede a todos; Litigante y Paralegal necesitan asignación vigente, y Paralegal no puede sellar. Cliente puede leer los metadatos del expediente asignado, pero conserva denegadas todas las operaciones documentales. Carga, sellado, verificación y exportación comprueban el expediente explícito y el documento; cambiar el expediente de la URL no permite acceder al documento, incluso como Owner. Las rutas documentales globales retiradas responden \texttt{404} y no invocan el workflow. Las pruebas también rechazan cabeceras ambiguas, campos que intentan suplantar actor o rol y cuerpos que exceden sus límites.

Retirar una asignación impide la siguiente operación con la misma sesión. Las pruebas separan dos órdenes concurrentes: una revocación confirmada impide el commit posterior; si la operación ya obtuvo los bloqueos que protegen la autorización, puede confirmar primero y la revocación espera. Asimismo, la revocación de sesión durante la preparación impide publicar el resultado. Una respuesta cuya autorización y auditoría se confirmaron antes de revocar puede terminar de transmitirse después; el ensayo no promete cancelar bytes ya autorizados.

\subsection{Atomicidad y sellado concurrente}\label{sec:pruebas-atomicidad}

Las pruebas de PostgreSQL inyectan fallos al insertar auditoría y al confirmar la transacción. Verifican que no quede un documento, asignación, usuario o consumo de recuperación confirmado sin su evento correspondiente. Escritores de documentos, expedientes e identidad extienden una sola cadena con secuencias únicas. La preparación criptográfica ocurre fuera de la transacción; el commit vuelve a comprobar identidad activa, rol, pertenencia y estado documental.

El ensayo de \rutarepo{scripts/api-concurrency-demo.sh}, invocado por la demostración HTTP, enfrenta dos procesos reales del servidor al sellado de una misma versión. Uno obtiene \texttt{200} y otro \texttt{409}; PostgreSQL conserva un solo evento de sellado y ambos servidores exportan un ZIP idéntico. La firma, el certificado con su CRL y el sello del paquete ganador se verifican con OpenSSL. Se prueba así la preservación de la evidencia ganadora, sin afirmar que dos preparaciones concurrentes realicen una única solicitud a la TSA.

La revisión de privilegios prueba roles operativos que pueden ejecutar el flujo permitido y rechaza capacidad de reescribir auditoría, modificar campos documentales inmutables, alterar recibos o asumir roles con esos permisos. Incluye propiedad de tablas y de la función protectora, privilegios por columna y esquemas modificables en la ruta de búsqueda. Estas pruebas no eliminan la confianza depositada en los administradores de PostgreSQL.

\subsection{Migración y restauración}\label{sec:pruebas-migracion}

El ensayo \rutarepo{scripts/api-migration-demo.sh} reconstruye un fixture de archivos documentales legacy a partir del cifrado y la evidencia reales de cuatro documentos de la demostración, tres de ellos sellados. La inspección comprueba un mapa completo hacia expedientes existentes, los bytes cifrados y la evidencia capturada. La aplicación conserva exactamente las 46 entradas históricas, incluidas sus marcas de tiempo y hashes, y añade su propio evento de importación. Un segundo intento no duplica documentos ni eventos. El respaldo con \texttt{pg\_dump} se restaura con \texttt{pg\_restore} en otra base: se comparan los datos SQL y los ZIP byte por byte, se reinicia el servidor con el rol restringido y se repite con OpenSSL la verificación independiente de firma, certificado con su CRL y sello.

El destino de esa importación se prepara con baselines explícitas: el guion copia los hechos originales de las raíces en un destino vacío y preserva el origen. Esa reconstrucción no demuestra la conversión íntegra del historial administrativo actual. El respaldo y la restauración completos del mismo guion comparan, además, todas las tablas, revisiones, autores y fechas, así como el detalle y la historia por HTTP.

Las regresiones cubren asociaciones ausentes, fuentes o mapas modificados tras inspeccionar, fallos de auditoría, barreras durables previas al commit, errores al escribir marcadores posteriores al commit y reintentos de reconciliación. El arranque rechaza una restauración parcial aunque conserve el recibo, así como cambios en documentos, asociaciones, prefijo de auditoría o marcadores. El origen se conserva para comparación y los adaptadores nuevos bloquean sus escrituras; las versiones antiguas del ejecutable deben permanecer detenidas durante el corte de archivos.

La validación histórica de cadena y CRL del firmante usa el instante del sello. La comprobación de la autoridad TSA conserva la política actual de OpenSSL y puede rechazar una autoridad expirada en la fecha de ejecución. Redis no comparte transacción con PostgreSQL: ante un fallo de auditoría, identidad intenta retirar desafíos o sesiones recién creados y no entrega sus tokens; si también falla esa limpieza, su caducidad limita la permanencia. Logout conserva la revocación aunque falle la auditoría posterior.
